\documentclass[12pt, a4paper]{article}

% ---- Standard Packages ----
\usepackage[a4paper, margin=1in]{geometry}
\usepackage[utf8]{inputenc}
\usepackage[T1]{fontenc}
\usepackage[Indonesian]{babel}
\usepackage{amsmath, amssymb, amsthm, mathtools}
\usepackage{enumitem}
\usepackage{booktabs}
\usepackage{multicol}
\usepackage{hyperref}

% ---- Basic Metadata ----
\title{MODUL BELAJAR SEJARAH KELAS 10\\ \large Persiapan Ulangan Akhir Semester}
\author{Ringkasan Materi Komprehensif}
\date{}

\begin{document}
\maketitle

\tableofcontents
\newpage

\section{Sejarah Secara Umum}
\subsection{Pengertian Sejarah}
Secara etimologis, kata sejarah memiliki keterkaitan dengan asal-usul kata dari berbagai bahasa asing yang mencerminkan makna pertumbuhan, silsilah, dan rekonstruksi masa lalu:
\begin{itemize}
	\item \textbf{Bahasa Arab (Sajaratun):} Berarti pohon. Kata ini merujuk pada silsilah, struktur keluarga, atau perkembangan suatu peristiwa yang tumbuh bercabang menyerupai pohon dari akar hingga ke ranting.
	\item \textbf{Bahasa Inggris (History):} Berasal dari kata yang merujuk pada masa lampau umat manusia yang disusun berdasarkan bukti-bukti tertulis maupun lisan.
	\item \textbf{Bahasa Yunani (Istoria):} Berarti ilmu atau orang pandai. Istoria bermakna sebagai hasil penelitian atau penyelidikan mendalam mengenai peristiwa manusia di masa lampau.
\end{itemize}

\subsubsection{Pengertian Sejarah Menurut Para Ahli}
\begin{itemize}
	\item \textbf{Herodotus (Bapak Sejarah):} Sejarah tidak berkembang ke depan dengan tujuan pasti, melainkan bergerak layaknya garis lingkaran yang tinggi rendahnya diakibatkan oleh keadaan perputaran jatuh bangunnya masyarakat dan peradaban manusia.
	\item \textbf{Aristoteles:} Sebuah sistem yang meneliti kejadian awal dan tersusun dalam bentuk kronologi peristiwa di masa lalu yang mempunyai catatan rekaman bukti konkret serta empiris.
	\item \textbf{Ibnu Khaldun:} Sejarah adalah catatan mengenai manusia dan peradabannya, yang mencakup proses perubahan sosial, watak masyarakat, serta sebab-akibat dari peristiwa-peristiwa masa lalu secara mendalam dan filosofis.
	\item \textbf{Sartono Kartodirdjo:} Sejarah merupakan rekonstruksi masa lampau dalam bentuk cerita yang disusun kembali secara verbal berdasarkan fakta-fakta sejarah yang dapat dipertanggungjawabkan.
\end{itemize}

\subsubsection{Sejarah Objektif}
Merujuk pada kejadian atau peristiwa itu sendiri yang benar-benar terjadi di masa lampau tanpa adanya pengaruh dari opini, pandangan personal, atau prasangka subjektif peneliti. Sejarah objektif tidak dapat diulang dan didukung oleh fakta keras.

\subsubsection{Sejarah Subjektif}
Merupakan produk rekonstruksi atau peristiwa sejarah yang dibangun oleh cerita atau hasil tulisan sejarawan. Sifatnya subjektif karena sudah melibatkan interpretasi, sudut pandang, emosi, dan latar belakang pemikiran dari penulis sejarah tersebut.

\subsection{Ruang Lingkup Sejarah}
Sejarah memiliki cakupan yang luas dan dapat dilihat dari empat sudut pandang utama:
\begin{enumerate}
	\item \textbf{Sejarah sebagai Peristiwa:} Kejadian nyata di masa lalu yang benar-benar terjadi, bersifat unik (hanya terjadi sekali), abadi (tidak berubah), dan penting (memberikan pengaruh besar bagi kehidupan masyarakat luas).
	\item \textbf{Sejarah sebagai Ilmu:} Suatu susunan pengetahuan (*body of knowledge*) tentang masa lalu yang disusun secara sistematis menggunakan metode ilmiah, memiliki objek, teori, serta generalisasi tertentu.
	\item \textbf{Sejarah sebagai Kisah:} Rekonstruksi peristiwa masa lalu yang diceritakan kembali oleh seseorang, baik secara lisan maupun tulisan melalui buku, majalah, atau media lainnya.
	\item \textbf{Sejarah sebagai Seni:} Penulisan sejarah yang membutuhkan keterlibatan intuisi, imajinasi, emosi, dan gaya bahasa yang indah agar pembaca dapat merasakan atmosfer peristiwa masa lalu tersebut secara hidup.
\end{enumerate}

\subsection{Pembagian Sejarah Berdasarkan Tema}
\begin{enumerate}
	\item \textbf{Sejarah Politik:} Fokus pada perkembangan kekuasaan, pemerintahan, kepemimpinan, birokrasi, konflik kekuasaan, diplomasi, dan kebijakan negara.
	\item \textbf{Sejarah Perekonomian:} Meneliti sistem produksi, distribusi, konsumsi, perdagangan, mata uang, perkembangan pasar, serta kondisi finansial masyarakat dari waktu ke waktu.
	\item \textbf{Sejarah Agama:} Mengkaji asal-usul, penyebaran, perkembangan pemikiran keagamaan, institusi suci, serta pengaruh nilai-nilai spiritual terhadap peradaban manusia.
\end{enumerate}

\subsection{Pembagian Sejarah Berdasarkan Periode Waktu (Periodisasi)}
Periodisasi dibuat untuk memudahkan pemahaman terhadap babakan waktu sejarah Indonesia:
\begin{enumerate}
	\item Zaman Praaksara (Sebelum mengenal tulisan)
	\item Zaman Hindu - Buddha (Masuknya pengaruh kebudayaan India)
	\item Zaman Islam (Penyebaran Islam dan berdirinya kesultanan)
	\item Zaman Penjajahan Barat (Kolonialisme Portugis, VOC, Belanda, dan Inggris)
	\item Zaman Pergerakan Nasional (Kebangkitan organisasi modern menuju kemerdekaan)
	\item Zaman Pendudukan Jepang (Fasisme militer Jepang di Nusantara)
	\item Zaman Proklamasi dan Kemerdekaan (1945)
	\item Zaman Revolusi Fisik (Mempertahankan kemerdekaan dari agresi Sekutu dan Belanda)
	\item Zaman Orde Lama (Pemerintahan di bawah Presiden Soekarno)
	\item Zaman Orde Baru (Pemerintahan di bawah Presiden Soeharto)
	\item Zaman Reformasi (1998 hingga masa kini)
\end{enumerate}

\subsection{Pembagian Sejarah dalam Ruang dan Waktu}
Sejarah selalu terikat oleh dua dimensi utama, yaitu ruang (*spasial*) dan waktu (*temporal*). Unsur waktu membatasi kapan peristiwa itu terjadi secara kronologis, sedangkan pembagian sejarah secara ruang membatasi cakupan geografis tempat peristiwa berlangsung menjadi tingkat lokal, regional wilayah, nasional, hingga global.

\subsubsection{Sejarah Sebagai Kisah}
Merupakan kejadian yang terjadi di masa lalu kemudian dibangun kembali lewat ingatan atau penulisan sejarah. Untuk menyusun kisah ini, diperlukan bukti sejarah yang valid:
\begin{enumerate}
	\item \textbf{Lisan (*Oral History*):} Keterangan langsung dari pelaku atau saksi sejarah yang masih hidup mengenai peristiwa yang mereka alami sendiri.
	\item \textbf{Tradisi Lisan (*Oral Tradition*):} Cerita rakyat, legenda, mitos, atau dongeng yang diwariskan secara turun-temurun dari satu generasi ke generasi berikutnya.
\end{enumerate}

\subsubsection{Sejarah Sebagai Ilmu}
Menurut \textit{E.H. Carr}, sejarah adalah proses interaksi yang berkesinambungan antara sejarawan dengan fakta-faktanya, sebuah dialog tanpa akhir antara masa kini dan masa lalu. Sejarah bukan sekadar kumpulan fakta mati, melainkan interpretasi aktif yang melibatkan konteks zaman sejarawan itu sendiri.

\noindent \textbf{Sifat Sejarah Sebagai Ilmu:}
\begin{enumerate}
	\item \textbf{Empiris:} Sejarah didasarkan pada pengalaman, pengamatan, dan bukti nyata dari aktivitas manusia di masa lampau, bukan rekayasa atau mitos.
	\item \textbf{Memiliki Objek:} Objek kajian sejarah adalah manusia dan waktu. Sejarah memiliki objek sendiri yang unik dan tidak dimiliki oleh penelitian ilmu lain, yaitu aktivitas manusia dalam dimensi waktu lampau (Contoh: Meneliti saluran masuknya Islam pada abad ke-7 hingga abad ke-13 Masehi).
	\item \textbf{Memiliki Teori:} Memiliki kumpulan kaidah, asas, dan prinsip ilmiah yang melandasi jalannya proses penelitian sejarah.
	\item \textbf{Memiliki Metode:} Memiliki langkah-langkah tersendiri untuk menjelaskan perkembangan, kesinambungan, atau perubahan peristiwa sejarah secara valid.
	\item \textbf{Generalisasi:} Sejarah mampu menarik kesimpulan umum dari fakta-fakta yang ditemukan, sekaligus bertindak sebagai alat koreksi atau komplementer atas kumpulan teori ilmu sosial lainnya.
\end{enumerate}

\subsubsection{Sejarah Sebagai Seni}
Dalam menulis sejarah, sejarawan memerlukan elemen seni agar tulisan tidak terasa kaku dan membosankan. Elemen tersebut meliputi:
\begin{enumerate}
	\item \textbf{Intuisi:} Pemahaman langsung atau insting sejarawan saat memilih fakta dan menentukan arah penulisan setelah memeriksa sumber daya.
	\item \textbf{Imajinasi:} Kemampuan sejarawan membayangkan bagaimana suasana, detail, dan kronologi peristiwa masa lampau itu terjadi berdasarkan data yang ada.
	\item \textbf{Emosi:} Keterlibatan perasaan sejarawan untuk berempati dengan objek penelitian tanpa kehilangan objektivitas ilmiahnya.
	\item \textbf{Gaya Bahasa:} Penggunaan tata bahasa yang menarik, komunikatif, hidup, dan mudah dipahami oleh pembaca.
\end{enumerate}

\section{Tahapan Penelitian Sejarah}
Penelitian sejarah adalah sebuah penelitian yang membahas tentang masa lalu manusia. Tujuan dari penelitian sejarah adalah untuk memahami, merekonstruksi, dan memecahkan kebenaran peristiwa di masa lalu secara kritis, baik mengenai penyebab (*kausalitas*), pengaruh, perkembangan, maupun akibat yang ditimbulkan dari suatu peristiwa.

\subsection{Langkah-Langkah Metode Sejarah}
\subsubsection{1. Pemilihan Topik}
Langkah awal di mana peneliti menentukan masalah atau tema yang akan diteliti. Pemilihan topik harus memenuhi dua syarat kedekatan:
\begin{itemize}
	\item \textbf{Kedekatan Emosional:} Kedekatan emosional menjadi bekal utama dan sangat dibutuhkan untuk menumbuhkan rasa suka, ketertarikan, dan motivasi yang memudahkan peneliti mengkaji suatu objek penelitian dalam jangka panjang.
	\item \textbf{Kedekatan Intelektual:} Seorang peneliti harus menguasai topik apa yang akan diteliti dengan mencari tahu lebih dalam mengenai teori, konsep, dan ketersediaan sumber daya dari topik yang hendak diteliti tersebut.
\end{itemize}

\subsubsection{2. Heuristik (Pengumpulan Sumber)}
Heuristik adalah tahapan pengumpulan data atau berburu sumber-sumber sejarah. Tahap ini dilakukan setelah peneliti menentukan topik penelitian. Keberhasilan seorang peneliti dalam mencari sumber ditentukan oleh wawasan dari peneliti atau sejarawan mengenai jenis, klasifikasi, serta lokasi penyimpanan sumber-sumber sejarah yang diperlukan (seperti arsip, perpustakaan, atau wawancara lapangan).

\subsubsection{3. Kritik Sumber (Verifikasi)}
Proses ini dilakukan untuk membuktikan bahwa sumber sejarah yang bersangkutan adalah asli (*otentisitas*) dan untuk menguji tingkat kepercayaan/kredibilitas dari isi sumber sejarah tersebut.
\begin{itemize}
	\item \textbf{Kritik Eksternal:} Menguji keaslian fisik sumber sejarah (misal: usia kertas, jenis tinta, atau bahan prasasti).
	\item \textbf{Kritik Internal:} Menguji kredibilitas isi informasi dari sumber sejarah, serta memastikan peran saksi/penulis dalam peristiwa sejarah yang sedang diteliti.
\end{itemize}

\subsubsection{4. Interpretasi}
Interpretasi adalah penafsiran makna dari fakta-fakta sejarah yang telah lolos uji verifikasi. Interpretasi dibagi menjadi dua proses utama:
\begin{enumerate}
	\item \textbf{Analisis:} Proses mengurai fakta sejarah untuk memperoleh penjelasan dari sumber sejarah yang tidak secara implisit atau gamblang menjelaskan suatu kejadian atau peristiwa.
	\item \textbf{Sintesis:} Proses menyatukan, menghubungkan, dan merangkai hasil analisis-analisis dari berbagai sumber sejarah guna mencapai benang merah dan kesimpulan utama yang utuh.
\end{enumerate}

\subsubsection{5. Historiografi}
Tahapan terakhir dalam penelitian sejarah adalah historiografi, yaitu proses penulisan sejarah. Fakta-fakta sejarah yang didapat dirangkai secara kronologis, kausalitas, dan sistematis menjadi satu kesatuan cerita atau karya ilmiah sejarah yang utuh.

\section{Pendekatan Penelitian dan Ilmu Bantu Sejarah}
\subsection{Pendekatan Penelitian}
Pendekatan dalam setiap penelitian sosial bermacam-macam, tetapi ilmu sejarah memiliki dua pendekatan khas yang membedakannya dengan ilmu sosial lain:
\begin{itemize}
	\item \textbf{Pendekatan Diakronis:} Memanjang dalam waktu tetapi menyempit dalam ruang. Fokus melihat perubahan kronologis dari waktu ke waktu (bersifat dinamis).
	\item \textbf{Pendekatan Sinkronis:} Meluas dalam ruang tetapi terbatas/menyempit dalam waktu. Fokus menganalisis struktur, fungsi, dan kondisi sosial-ekonomi mendalam pada suatu masa tertentu (bersifat statis dan struktural).
\end{itemize}

\subsection{Ilmu Bantu Sejarah}
Ilmu bantu sejarah adalah disiplin ilmu lain yang berkaitan dengan kehidupan manusia sehari-hari dan digunakan untuk membantu sejarawan menganalisis sumber sejarah secara akurat:
\subsubsection{Arkeologi}
Disebut juga ilmu purbakala. Berkaitan dengan bekas atau warisan masa lalu berupa benda fisik purba atau artefak (benda visual hasil buatan manusia).
\subsubsection{Paleontologi}
Ilmu yang mengkaji bentuk-bentuk kehidupan purba, perkembangan flora, dan fauna kuno yang pernah ada di bumi berdasarkan sisa-sisa organik, terutama fosil.
\subsubsection{Paleoantropologi}
Merupakan gabungan dari dua disiplin ilmu yaitu paleontologi dan antropologi. Objek kajian paleoantropologi adalah mempelajari asal-usul, evolusi fisik, ras, dan kebudayaan dari fosil manusia purba.
\subsubsection{Paleografi}
Digunakan untuk mengkaji manuskrip, aksara, atau tulisan-tulisan kuno yang terdapat pada dokumen kertas, kulit, maupun daun lontar sebagai satu sumber sejarah primer.
\subsubsection{Epigrafi}
Cabang ilmu yang mempelajari tulisan kuno yang dipahat pada material keras, seperti batu (prasasti) atau logam kuno.
\subsubsection{Ikonografi}
Ilmu yang mempelajari tentang arca, patung-patung kuno, simbol, atau relief keagamaan masa lampau.
\subsubsection{Numismatik}
Ilmu yang mengkaji tentang mata uang kuno, medalion, atau segel perdagangan. Numismatik memberikan bahan berharga bagi sejarah ekonomi dan seni.
\subsubsection{Keramologi}
Ilmu yang mempelajari tentang benda bersejarah berupa keramik, tembikar, wadah tanah liat, pecahan porselen, dan menentukan asal serta usia pembuatannya.
\subsubsection{Genealogi}
Merupakan cabang ilmu yang mempelajari asal-usul silsilah keturunan, hubungan keluarga, dan nenek moyang suatu tokoh atau bangsa.
\subsubsection{Filologi}
Ilmu yang berkaitan dengan bahasa, teks-teks sastra kuno, naskah manuskrip lama, dan perubahan struktur kebahasaan dari masa ke masa.

\section{Konsep Berpikir Sejarah}
Konsep berpikir sejarah dibagi menjadi dua pilar utama, yaitu:
\begin{enumerate}
	\item Berpikir diakronis (kronologis, menekankan urutan waktu)
	\item Berpikir sinkronis (struktural, menekankan analisis sistem)
\end{enumerate}

\subsection{Cara Berpikir Sinkronis}
Sinkronis artinya segala sesuatu yang bersangkutan dengan peristiwa yang terjadi pada suatu media atau ruang dalam waktu yang terbatas. Pendekatan ini membedah gejala sosial secara mendalam (misal: Keadaan ekonomi masyarakat Jawa pada masa Tanam Paksa tahun 1830).

\subsection{Kehidupan Manusia, Ruang, dan Waktu}
Tiga unsur penting dalam sejarah yang tidak bisa dipisahkan satu sama lain:
\subsubsection{Manusia}
Manusia adalah aktor utama dalam sejarah. Sejarah adalah ilmu tentang manusia.
\begin{itemize}
	\item \textbf{Objek:} Manusia sebagai aktor yang melakukan dan membentuk setiap kejadian sejarah.
	\item \textbf{Subjek:} Manusia berperan aktif merekonstruksi dan menulis kembali sebuah kejadian masa lalu.
\end{itemize}
\noindent Filsof Yunani \textit{Heraclitus} menyatakan: ``\textit{Panta rhei uden menei}'', yang berarti ``Tidak ada yang tetap, semua mengalir karena mengalami perkembangan dan perubahan.''

\subsubsection{Ruang atau Spasial}
Ruang merupakan suatu tempat geografis atau wadah terjadinya peristiwa alam, sosial, serta peristiwa sejarah.
\begin{itemize}
	\item \textbf{Konsep Sejarah - Ruang - Penelaahan:} Suatu peristiwa tidak dapat dipisahkan dari dimensi ruang dan waktu. Ruang memberikan konteks geografis di mana manusia bertindak, sementara waktu menitikberatkan pada aspek ``kapan'' peristiwa itu terjadi.
\end{itemize}

\subsubsection{Waktu}
Menurut sejarawan Kuntowijoyo, konsep waktu dalam sejarah mencakup 4 dimensi utama:
\begin{enumerate}
	\item \textbf{Perkembangan:} Terjadi jika masyarakat mengalami pergerakan berturut-turut dari satu bentuk ke bentuk yang lain yang lebih kompleks.
	\item \textbf{Kesinambungan:} Terjadi jika masyarakat baru tetap mengadopsi atau melanjutkan lembaga, kebiasaan, atau tradisi dari masa lalu.
	\item \textbf{Pengulangan:} Terjadi jika peristiwa di masa lalu kembali terjadi di masa sesudahnya dengan pola atau karakteristik yang mirip, meskipun pelaku dan waktunya berbeda.
	\item \textbf{Perubahan:} Terjadi jika masyarakat mengalami pergeseran secara besar-besaran dan dalam waktu yang relatif singkat, biasanya didorong oleh faktor eksternal.
\end{enumerate}

\subsubsection{Perbedaan Sinkronis dan Diakronis} 
\begin{enumerate}
	\item \textbf{Ciri-ciri Diakronis:}
		\begin{enumerate}
			\item Mengkaji peristiwa seiring berjalannya waktu (bersifat vertikal).
			\item Menitikberatkan kepada proses perkembangan dan kronologi peristiwa sejarah.
			\text{Cakupan waktu luas, namun cakupan ruang pembahasan sempit.}
		\end{enumerate}
	\item \textbf{Ciri-ciri Sinkronis:}
		\begin{enumerate}
			\item Mengkaji masa tertentu secara mendalam (bersifat horizontal).
			\item Menitikberatkan kepada kultur, sistem sosial, ekonomi, politik, dan karakter struktural.
			\item Cakupan waktu sempit, namun cakupan ruang atau variabel bahasan sangat luas dan mendalam.
		\end{enumerate}
\end{enumerate}

\section{Sumber Sejarah}
Eksistensi manusia pada masa lalu dapat diketahui melalui berbagai peninggalannya. Peninggalan peradaban masa lalu dapat berupa manuskrip kuno, bangunan megah, maupun benda-benda historis lainnya. Jejak peristiwa tersebut ada yang berbentuk lisan, tulisan, maupun benda. Jejak-jejak sejarah inilah yang kemudian digunakan sejarawan sebagai dasar atau instrumen untuk merekonstruksi peristiwa sejarah sehingga dapat menemukan kebenaran masa lalu.

\subsection{Sumber Sejarah Berdasarkan Bentuk}
\begin{itemize}
	\item \textbf{Sumber Tertulis:} Sumber sejarah yang diperoleh dari benda peninggalan yang memuat aksara/tulisan (Contoh: kitab, prasasti, dokumen kearsipan, babad, kronik, surat kabar kuno).
	\item \textbf{Sumber Lisan:} Sumber sejarah yang diperoleh melalui wawancara langsung (*oral history*) atau penuturan lisan dari pelaku dan saksi sejarah, atau berupa tradisi lisan (*oral tradition*).
	\item \textbf{Sumber Audio-Visual:} Sumber sejarah yang berbentuk kombinasi rekaman suara dan gambar bergerak (Contoh: film dokumenter sejarah, rekaman video berita zaman lampau, kaset pidato).
\end{itemize}

\subsection{Sumber Sejarah Berdasarkan Sifat (Kedudukan)}
\begin{itemize}
	\item \textbf{Sumber Primer:} Sumber yang sifatnya asli, sezaman, dan berasal langsung dari pelaku atau saksi mata yang berada di lokasi kejadian (Contoh: prasasti, arsip resmi negara, saksi hidup, benteng peninggalan).
	\item \textbf{Sumber Sekunder:} Sumber sejarah yang berasal dari pihak kedua yang tidak mengalami peristiwa secara langsung, melainkan hasil kajian dari peneliti atau sejarawan berdasarkan sumber primer (Contoh: buku teks sejarah, biografi, ensiklopedia).
\end{itemize}
\noindent \textbf{Bukti Sejarah:} Sesuatu yang memperkuat dan membenarkan rekonstruksi suatu sejarah.\\
\noindent \textbf{Fakta Sejarah:} Kesimpulan akhir atau hasil analisis kritis yang diperoleh dari pengolahan berbagai bukti-bukti sejarah.

\subsection{Sumber Sejarah Berdasarkan Wujudnya}
\begin{itemize}
	\item \textbf{Artefak:} Benda fisik berupa perkakas atau alat-alat yang dibuat dan digunakan oleh manusia masa lalu (Contoh: kapak perimbas, mata panah, perhiasan logam, arca).
	\item \textbf{Fakta Mental (Mentofak):} Kondisi psikologis, keyakinan, ideologi, alam pikiran, dan tatanan nilai mental yang berkembang dalam masyarakat pada periode tertentu.
	\item \textbf{Fakta Sosial (Sosiofak):} Perkembangan struktur kondisi sosial, sistem kemasyarakatan, jaringan interaksi, hukum adat, dan perilaku kolektif antarmanusia.
\end{itemize}

\subsection{Fakta Sejarah Berdasarkan Sifat Kebenarannya}
\begin{itemize}
	\item \textbf{Fakta Keras (*Hard Fact*):} Fakta sejarah yang kebenarannya sudah mutlak, pasti, teruji, dan tidak dapat diperdebatkan lagi oleh siapa pun (Contoh: Proklamasi Kemerdekaan Indonesia terjadi pada 17 Agustus 1945).
	\item \textbf{Fakta Lunak (*Soft Fact*):} Fakta sejarah yang kebenarannya masih belum jelas, bersifat hipotesis, dan masih terbuka untuk diperdebatkan atau diteliti lebih lanjut karena minimnya bukti pendukung (Contoh: kepastian letak pusat pusat Kerajaan Sriwijaya).
\end{itemize}

\section{Konsep Hidup Manusia dalam Perubahan dan Berkelanjutan}
\subsection{Perubahan}
Merupakan suatu fenomena di mana masyarakat mengalami pergeseran dari satu keadaan ke keadaan lainnya yang berbeda dari waktu ke waktu secara nyata, baik dalam aspek sosial, budaya, ekonomi, maupun politik.

\subsection{Keberlanjutan (Kontinuitas)}
Merupakan rangkaian peristiwa yang telah terjadi di masa lampau, yang terus diadaptasi atau dilanjutkan pengaruhnya hingga masa sekarang dan masa yang akan datang tanpa adanya interupsi total.

\subsubsection{Cakupan Sejarah}
Ilmu sejarah mencakup tiga dimensi waktu yang saling bertautan erat:
\begin{enumerate}
	\item Penglihatan masa lalu (akar fondasi kehidupan)
	\item Masa sekarang (konteks aksi manusia saat ini)
	\item Masa yang akan datang (pembelajaran untuk proyeksi masa depan)
\end{enumerate}

\section{Sejarah Bagi Kehidupan Manusia}
Sejarah memiliki peran dan arti penting dalam kehidupan bermasyarakat dan bernegara.
\subsection{Fungsi Sejarah Bagi Kehidupan Manusia}
Secara garis besar sejarah memiliki fungsi genetis (menjelaskan asal-usul), didaktis (memberikan edukasi dan pelajaran berharga), instruktif (sebagai alat bantu mengajar), serta inspiratif. Sejarah membantu manusia mengambil hikmah dari kekeliruan masa lalu dan menjadikannya teladan emas bagi masa depan.
\subsection{Fungsi Sejarah Menurut \textit{E.H. Carr}}
Dapat memuaskan rasa ingin tahu tentang kehidupan kolektif orang lain, serta berfungsi sebagai jembatan komparatif yang membandingkan kehidupan masa sekarang dengan dinamika masa lalu guna menemukan solusi peradaban.

\section{Jenis Manusia Purba di Indonesia}
\subsection{Meganthropus paleojavanicus}
Ditemukan oleh \textit{Gustav Heinrich Ralph von Koenigswald} di daerah Sangiran pada tahun 1936 dan 1941. Manusia purba ini hidup pada masa pleistosen awal.
\noindent \textbf{Ciri-Ciri:}
\begin{itemize}
	\item Berbadan besar, tegap, dan memiliki tonjolan tajam di bagian belakang kepala.
	\item Memiliki tulang pipi yang tebal dan kuat.
	\item Tonjolan kening sangat mencolok dan tebal.
	\item Tidak memiliki dagu.
	\item Memiliki otot rahang dan gigi geligi yang sangat kuat untuk mengunyah tumbuh-tumbuhan.
\end{itemize}

\subsection{Pithecanthropus}
Jenis manusia purba dengan kerangka tubuh yang lebih muda dan paling banyak ditemukan di Indonesia. Berarti ``manusia kera yang berjalan tegak.''
\subsubsection{Pithecanthropus mojokertensis}
Pithecanthropus mojokertensis adalah manusia jenis Pithecanthropus yang paling tua. Fosilnya ditemukan di Kepuhklagen, Mojokerto, Jawa Timur pada tahun 1936 oleh Von Koenigswald.
\noindent \textbf{Ciri-Ciri:}
\begin{itemize}
	\item Badan tegap, tinggi badan berkisar antara 165--180 cm.
	\item Tulang rahang atas bawah sangat kuat dan otot pengunyah kokoh.
	\item Bagian kening menonjol tebal melintang.
	\item Hidung lebar, tulang pipi kuat, dan tidak memiliki dagu.
	\item Volume otak belum sempurna, kapasitasnya berkisar antara 750 - 1300 cc.
\end{itemize}

\subsubsection{Pithecanthropus erectus}
Berarti manusia kera yang berjalan tegak. Hidup sekitar 2,5 juta - 1,5 juta tahun yang lalu. Ditemukan oleh \textit{Eugene Dubois} di Trinil, Ngawi, Jawa Timur pada tahun 1891.
\noindent \textbf{Ciri-Ciri:}
\begin{itemize}
	\item Volume otak berkisar antara 750 - 1350 cc.
	\item Tinggi badan sekitar 160 - 180 cm.
	\item Postur tubuh tegap, namun tidak setegap Meganthropus.
	\item Gigi geraham yang besar dan rahang yang sangat kuat.
	\item Hidung tebal dan lebar.
	\item Tonjolan kening tebal melintang di dahi dari pelipis ke pelipis.
\end{itemize}

\subsubsection{Pithecanthropus soloensis}
Hidup pada era pleistosen atas (sekitar 900.000 - 300.000 tahun yang lalu). Ditemukan oleh Von Koenigswald dan Oppenoorth di daerah Ngandong dan Sambungmacan, dekat Bengawan Solo.
\noindent \textbf{Ciri-Ciri:}
\begin{itemize}
	\item Hidup kala pleistosen awal hingga pleistosen tengah.
	\item Tinggi badan berkisar antara 165 - 180 cm.
	\item Ketika berjalan, posisi tubuh sudah tegak lurus sempurna.
	\item Volume otak berkembang lebih baik, berkisar antara 1000 - 1300 cc.
	\item Alat pengunyah masih kuat namun bentuk wajah mengecil.
	\item Bagian kening menonjol tebal, hidung lebar, tidak memiliki dagu, bagian belakang tengkorak menonjol tajam.
\end{itemize}

\subsection{Homo}
Jenis manusia purba yang tingkat evolusinya paling tinggi dan hidup pada era pleistosen akhir serta memiliki kemiripan fisik dengan manusia modern.
\subsubsection{Homo wajakensis}
Ditemukan di Wajak, Tulungagung, Jawa Timur oleh \textit{B.D. van Rietschoten} pada tahun 1889 dan diteliti oleh Eugene Dubois.
\begin{itemize}
	\item Volume otak sudah maju, mencapai sekitar 1630 cc.
	\item Muka datar dan berukuran cukup lebar.
	\item Rahang tergolong padat, besar, dan memiliki gigi-gigi berukuran besar.
	\item Tinggi badan sekitar 173 cm, struktur tubuh tegak.
\end{itemize}
\subsubsection{Homo floresiensis}
Ditemukan di dalam gua Liang Bua, Pulau Flores, Nusa Tenggara Timur oleh tim arkeologi gabungan Indonesia dan Australia pada tahun 2003. Sering dijuluki sebagai ``hobbit.''
\begin{itemize}
	\item Tinggi badan kerdil, hanya mencapai sekitar 1 meter (100 cm).
	\item Bentuk dahi sempit, rata, dan tidak menonjol.
	\item Kerangka tulang kepala dan tengkorak berukuran sangat kecil.
	\item Volume otak primitif, hanya sekitar 380 cc.
	\item Berat badan rata-rata 25 kg.
	\item Tulang rahang menonjol ke depan namun tidak memiliki dagu modern.
\end{itemize}
\subsubsection{Homo soloensis}
Ditemukan pada tahun 1931 - 1933 oleh \textit{Ter Haar, W.F.F. Oppenoorth}, dan Von Koenigswald di Ngandong, lembah Bengawan Solo.
\begin{itemize}
	\item Volume otak maju, berkisar antara 1000 - 1300 cc.
	\item Tinggi badan dapat mencapai 180 - 210 cm.
	\item Struktur tulang wajah tidak semirip manusia purba awal, otot tengkuk menyusut.
\end{itemize}
\subsubsection{Homo sapiens}
Merupakan jenis manusia purba cerdas yang menjadi cikal bakal langsung manusia modern (*Anatomically Modern Human*).
\begin{itemize}
	\item Volume otak berkisar antara 1350 - 1450 cc.
	\item Tinggi badan bervariasi antara 130 - 210 cm.
	\item Berat badan antara 30 - 150 kg, memiliki dagu nyata, dan otot kunyah mereduksi.
\end{itemize}

\section{Corak Kehidupan Manusia Purba}
Prof. Dr. R. \textit{Soekmono} membagi pembabakan zaman praaksara secara arkeologis menjadi dua periodisasi besar:
\begin{enumerate}
	\item \textbf{Zaman Batu:} Terdiri atas Paleolitikum (Batu Tua), Mesolitikum (Batu Tengah), Neolitikum (Batu Muda), dan Megalitikum (Batu Besar).
	\item \textbf{Zaman Logam:} Terdiri atas Zaman Tembaga, Zaman Perunggu, dan Zaman Besi. Di Indonesia, Zaman Logam yang dialami secara dominan adalah Zaman Perunggu dan Zaman Besi (Zaman Perundagian).
\end{enumerate}

\subsection{Periodisasi Berdasarkan Pola Mata Pencaharian}
\begin{enumerate}
	\item Masa berburu dan mengumpulkan makanan tingkat awal (*Paleolitikum*)
	\item Masa berburu dan mengumpulkan makanan tingkat lanjut (*Mesolitikum*)
	\item Masa bercocok tanam (*Neolitikum* dan *Megalitikum*)
	\item Masa perundagian (*Zaman Logam*)
\end{enumerate}

\subsubsection{Zaman Paleolitikum (Zaman Batu Tua)}
Peralatan dan senjata pada zaman ini dibuat dari batu mentah, pengerjaannya masih sangat kasar dan belum diasah. Manusia yang hidup pada zaman ini adalah Meganthropus dan Pithecanthropus. 
Mereka hidup berkelompok kecil (10 - 15 orang) demi keamanan. Besarnya kelompok ditentukan oleh kelimpahan hasil perburuan di alam sekitar. Kelompok ini memisahkan diri dan berpindah tempat dengan cara migrasi jika sumber pangan habis. Kaum laki-laki bertugas mencari binatang buruan, sedangkan perempuan mengurus anak dan memungut buah liar.

Dua hal penting yang ditemukan pada masa berburu dan mengumpulkan makanan:
\begin{enumerate}
	\item Dikenalkannya penggunaan api untuk memasak dan mengusir binatang buas.
	\item Alat-alat batu kasar mulai digunakan sebagai perkakas pertahanan hidup.
\end{enumerate}
Mereka mengumpulkan makanan langsung dari pemberian alam tanpa mengolahnya (*food gathering*). Mereka belum mampu memproduksi makanan, oleh karena itu mereka hidup berpindah-pindah tempat (*nomaden*) mengikuti gerak migrasi hewan buruan dan ketersediaan air bersih. Pada masa ini, kebudayaan dibagi menjadi dua pusat perkembangan: Kebudayaan Pacitan (kapak genggam) dan Kebudayaan Ngandong (alat tulang dan flakes).

\subsubsection{Zaman Neolitikum (Masa Bercocok Tanam dan Beternak)}
Zaman Neolitikum atau batu muda merupakan fase revolusi kebudayaan terbesar (*Neolithic Revolution*) bagi kehidupan masyarakat praaksara. Ciri utama masa Neolitikum:
\begin{itemize}
	\item Mulai meninggalkan kultur nomaden dan beralih ke pola hidup menetap secara permanen (*sedenter*).
	\item Mengembangkan pertanian bercocok tanam dengan metode tebang-bakar lahan (*slash and burn*).
	\item Terjadi perubahan radikal dari pola mencari makanan (*food gathering*) ke pola memproduksi makanan sendiri (*food producing*).
\end{itemize}
Mereka tidak lagi menggantungkan nasib sepenuhnya pada alam liar. Cara hidup berkelompok mulai terorganisasi dengan baik (20 - 50 orang) membentuk sebuah perkampungan awal. Mereka melakukan domestikasi hewan ternak (seperti ayam, kerbau, anjing) dan tanaman (padi, umbi-umbian) sehingga variasi makanan melimpah. Hasil pertanian yang melimpah disimpan di lumbung atas rumah sebagai persediaan logistik hingga musim panen berikutnya tiba. Alat batu sudah diasah sangat halus, seperti kapak persegi dan kapak lonjong.

\subsubsection{Masa Perundagian (Masa Logam dan Megalitikum)}
Masa Megalitikum (Batu Besar) berjalan beriringan dengan masa perundagian, memunculkan bangunan-bangunan pemujaan berukuran besar yang mempengaruhi sistem kepercayaan nenek moyang. Selanjutnya, manusia praaksara mengalami evolusi peradaban tinggi dengan ditemukannya teknik melebur bijih logam.

Zaman perundagian didasari atas munculnya golongan ``Undagi'', yaitu sekelompok tenaga kerja ahli yang memiliki keahlian khusus dalam pembuatan peralatan dan perhiasan dari bahan tembaga, perunggu, dan besi. Ciri-ciri kehidupan sosial masyarakat pada masa perundagian adalah:
\begin{itemize}
	\item Menganut kepercayaan animisme (pemujaan roh nenek moyang) dan dinamisme (kepercayaan benda keramat).
	\item Bermukim tetap secara teratur di daerah pegunungan, dataran rendah, hingga pesisir pantai.
	\item Hidup dalam kelompok masyarakat adat yang besar, terstruktur, dan memiliki sistem pembagian kerja.
	\item Memiliki kemampuan teknologi tinggi dalam mengolah bijih logam dengan menggunakan teknik cetak batu (*bivalve*) maupun teknik cetak lilin (*a cire perdue*).
	\item Berkembangnya pengetahuan tentang astronomi untuk pelayaran dan penanggalan pertanian.
	\item Perdagangan sistem barter berkembang pesat antar-pulau menggunakan perahu bercadik.
	\item Terbentuk struktur sosial yang dipimpin oleh seorang kepala suku bergelar *Primus Inter Pares*.
\end{itemize}

\section{Hasil Kebudayaan Manusia Purba Di Indonesia}
\subsection{Zaman Paleolitikum}
Tiga macam tradisi pembuatan alat yang dominan pada masa ini, yaitu:
\begin{enumerate}
	\item Tradisi kapak perimbas - penetak (*chopper - chopping tool*)
	\item Tradisi serpih bilah (*Flakes*)
	\item Tradisi alat tulang dan tanduk rusa (*Ngandong Culture*)
\end{enumerate}
\subsubsection{Tradisi Kapak Perimbas-Penetak}
Dikenal luas dalam dunia arkeologi dengan nama ``Kebudayaan Pacitan''. Contoh peninggalan budaya Pacitan adalah: Kapak perimbas (*chopper*), kapak penetak (*chopping tool*), pahat genggam (*proto hand adze*), kapak genggam awal (*proto hand axe*), kapak genggam standar (*hand axe*), dan serut genggam (*scraper*).

\subsection{Zaman Mesolitikum (Zaman Batu Tengah)}
Peralihan zaman dari Paleolitikum ke Neolitikum. Manusia purba pada masa ini mulai menunjukkan kemajuan pola hidup menetap sementara. Zaman ini ditandai oleh tiga macam kebudayaan atau tradisi utama:
\begin{enumerate}
	\item Kebudayaan Serpih Bilah (*Toala Culture*)
	\item Kebudayaan Alat Tulang (*Sampung Bone Culture*)
	\item Kebudayaan Kapak Genggam Sumatra (*Pebble/Sumatralith*)
\end{enumerate}
\subsubsection{Tradisi Serpih Bilah (Kebudayaan Toala)}
Kebudayaan serpih bilah (*flake \& blade*) mendapat pengaruh kuat dari unsur mikrolit sehingga menghasilkan alat batu kecil tajam yang mirip dengan batu api purba di Eropa. Ciri khas kebudayaan Toala adalah penggunaan alat serpih bergigi tajam. Tradisi serpih bilah ini dibedakan menjadi: pisau batu, alat serut, lancipan mata panah untuk berburu, dan mikrolit. Peneliti membuktikan persebaran kebudayaan Toala berada di Jawa Timur, Sulawesi Selatan, hingga Flores (sekitar 1000 SM).
\subsubsection{Tradisi Alat Tulang (Sampung Bone Culture)}
Ditemukan pada tahun 1928 - 1931 oleh \textit{Van Stein Callenfels} di Gua Lawa, Sampung, Ponorogo, Jawa Timur. Di lokasi ini ditemukan banyak perkakas yang dibuat dari tulang dan tanduk rusa. Pada masa Mesolitikum ini, manusia purba tinggal menetap sementara di dalam gua-gua alam yang disebut \textit{Abris Sous Roche}. Di daerah pantai, mereka juga meninggalkan tumpukan sampah dapur berupa kulit kerang dan siput yang membatu, yang disebut \textit{Kjokkenmoddinger}.

\section{Proses Migrasi Ras ke Nusantara}
\subsection{Proto Melayu (Bangsa Melayu Tua atau Austronesia)}
Masuk ke Indonesia sekitar tahun 1500 SM. Berasal dari wilayah Yunnan di Cina Selatan, kemudian bergerak menyusuri jalur darat Indo-Cina menuju Semenanjung Malaya. Dari sana, mereka menyebar melalui jalur laut ke kepulauan Nusantara menjadi dua jalur: Jalur Barat (membawa kebudayaan kapak persegi) dan Jalur Timur melalui Filipina (membawa kebudayaan kapak lonjong) menuju Melanesia dan Polinesia. Keturunan Proto Melayu di Indonesia saat ini antara lain suku Dayak, Toraja, Nias, dan Batak.

\subsection{Deutro Melayu (Bangsa Melayu Muda)}
Masuk ke Indonesia pada gelombang kedua sekitar tahun 500 SM. Rute migrasinya diawali dari Yunnan (Teluk Tonkin), Vietnam (membawa Kebudayaan Dongson), menyusuri Malaysia, hingga masuk ke Nusantara melalu jalur barat. Bangsa Deutro Melayu hidup pada Zaman Logam. \textit{H.R. van Heekeren} menyebut zaman ini sebagai zaman perundagian (*The Bronze - Iron Age*). Keturunan Deutro Melayu sangat unggul di bidang pertanian sawah basah, irigasi, dan pelayaran perahu bercadik. Suku bangsanya saat ini meliputi suku Jawa, Bugis, Melayu, Minang, dan Bali.

\subsection{Melanesoid (Akar Ras Negroid)}
Kedatangan mereka terjadi pada akhir Zaman Es (Pleistosen Akhir). Populasi nenek moyang Austro-Melanesoid ini bermigrasi dari daratan Asia menuju timur Nusantara. Sisa-sisa fisik keturunan awal mereka memiliki kedekatan genetik dengan Homo Wajakensis. Keturunannya saat ini menghuni wilayah Papua, Alor, dan NTT.

\subsection{Negroid dan Weddid}
\begin{itemize}
	\item \textbf{Negroid:} Penduduk asli Nusantara berkulit gelap, rambut keriting, dan bertubuh kecil mirip suku Semang di Malaya dan suku Negrito di Filipina.
	\item \textbf{Weddid (Bangsa Weda/Vedda):} Berasal dari Sri Lanka dengan fisik bertubuh kecil, kulit cokelat, dan rambut berombak. Masuk pada gelombang awal es mencair. Peninggalan artefak mereka meliputi: perkakas Mesolitikum berupa flakes, scraper, mata panah di Gua Dadu Murni, kapak lonjong dan persegi di Danau Sentani (Neolitikum), altar batu (Megalitikum) di Pulau Adi, serta corong perunggu di Danau Aimaru Papua.
\end{itemize}

\section{Sejarah Kerajaan Hindu-Buddha di Nusantara}
\subsection{Kerajaan Kutai}
Kerajaan Kutai Martadipura merupakan kerajaan bercorak Hindu tertua di Nusantara yang berdiri sekitar abad ke-4 Masehi. Kerajaan ini terletak di kawasan Muara Kaman, tepat di tepi hulu Sungai Mahakam, Kalimantan Timur.
\begin{itemize}
	\item \textbf{Sumber Sejarah:} Keberadaan Kutai dibuktikan dengan penemuan 7 buah prasasti berbentuk tugu batu pemujaan bernama \textit{Yupa}. Prasasti ini dipahat menggunakan huruf Pallawa dan ditulis dalam bahasa Sanskerta kuno.
	\item \textbf{Silsilah Raja:} Raja pertama sekaligus pendiri adalah Kudungga (nama asli Nusantara, menunjukkan pengaruh Hindu belum masuk sepenuhnya). Kudungga digantikan oleh putranya, Aswawarman, yang dianggap sebagai \textit{Wangsakarta} (pendiri dinasti) setelah melakukan upacara \textit{Wratyastoma} (upacara penyucian diri untuk masuk kasta Hindu). Raja terbesar Kutai adalah Mulawarman, putra Aswawarman.
	\item \textbf{Puncak Kejayaan:} Di bawah pemerintahan Raja Mulawarman, Kutai mencapai masa keemasan. Ia dikenal sangat dermawan, dibuktikan dengan catatan pada \textit{Yupa} bahwa ia menyedekahkan 20.000 ekor sapi kepada kaum Brahmana di tanah suci \textit{Waprakeswara} (tempat pemujaan Dewa Siwa).
	\item \textbf{Kehidupan Ekonomi:} Berpusat pada sektor peternakan, pertanian pedalaman, dan kendali atas perdagangan jalur sungai Mahakam yang menghubungkan dunia internasional.
\end{itemize}

\subsection{Kerajaan Tarumanegara}
Kerajaan Tarumanegara adalah kerajaan bercorak Hindu tertua di Pulau Jawa yang berdiri sekitar abad ke-5 Masehi (400--600 M) dan berpusat di wilayah Jawa Barat (sekitar Bogor, Bekasi, Karawang, dan Jakarta).
\begin{itemize}
	\item \textbf{Sumber Sejarah:} Keberadaannya dibuktikan melalui 7 buah prasasti penting, antara lain:
		\begin{enumerate}
			\item \textbf{Prasasti Ciaruteun:} Terdapat cap sepasang telapak kaki Raja Purnawarman yang dianalogikan sebagai kaki Dewa Wisnu (pemujaan aliran Hindu Waisnawa).
			\item \textbf{Prasasti Tugu:} Mengisahkan pembangunan/penggalian Sungai Candrabhaga dan Gomati sepanjang 6.112 busur (sekitar 11 km) untuk irigasi pertanian dan pencegahan banjir.
			\item Prasasti Jambu, Kebon Kopi, Pasir Awi, Muara Cianten, dan Cidanghiang.
		\end{enumerate}
	\item \textbf{Raja Terkenal:} Raja Purnawarman, seorang pemimpin yang gagah berani, bijaksana, dan sangat memperhatikan kesejahteraan ekonomi rakyatnya.
	\item \textbf{Kehidupan Sosial-Keagamaan:} Masyarakat menganut agama Hindu aliran Waisnawa, hidup berdampingan dengan sisa tradisi lokal kebudayaan megalitikum. Sektor ekonomi bertumpu pada pertanian padi serta peternakan sapi berskala besar.
\end{itemize}

\subsection{Kerajaan Sriwijaya}
Kerajaan Sriwijaya merupakan kerajaan maritim bernuansa Buddha terbesar di Asia Tenggara yang berdiri sejak abad ke-7 hingga ke-14 Masehi, berpusat di Palembang, Sumatra Selatan.
\begin{itemize}
	\item \textbf{Sumber Sejarah:} Berasal dari prasasti lokal seperti Prasasti Kedukan Bukit (menceritakan perjalanan suci atau \textit{Siddhayatra} Dapunta Hyang dengan membawa 20.000 tentara), Prasasti Talang Tuo (pembangunan Taman Sriksetra), Prasasti Telaga Batu, Kota Kapur, dan Karang Brahi. Selain itu, terdapat catatan musafir Cina bernama I-Tsing yang menyebutkan Sriwijaya sebagai pusat pembelajaran agama Buddha terbesar di luar India.
	\item \textbf{Raja Terkenal:} Balaputradewa (abad ke-9 M), putra dari Samaragrawira (Dinasti Syailendra), membawa Sriwijaya ke puncak kejayaan politik dan ekonomi. Ia menjalin hubungan internasional erat dengan Kerajaan Nalanda di India kuno.
	\item \textbf{Peran Strategis:} Sebagai kerajaan maritim, Sriwijaya menguasai jalur perdagangan internasional vital di Selat Malaka, Selat Sunda, dan Laut Cina Selatan melalui kekuatan armada lautnya yang kuat. Sriwijaya juga bertindak sebagai pusat kajian ilmiah agama Buddha aliran Vajrayana di Asia.
\end{itemize}

\subsection{Kerajaan Majapahit}
Kerajaan Majapahit merupakan kerajaan Hindu-Buddha terbesar di Nusantara yang berdiri pada tahun 1293 M hingga runtuh pada awal abad ke-16 M. Kerajaan ini berpusat di wilayah Mojokerto dan Trowulan, Jawa Timur.
\begin{itemize}
	\item \textbf{Pendirian:} Didirikan oleh Raden Wijaya yang berhasil memanfaatkan kedatangan pasukan Mongol (Dinasti Yuan) untuk menghancurkan Jayakatwang dari Kediri, sebelum akhirnya berbalik mengusir pasukan Mongol tersebut. Raden Wijaya dinobatkan bergelar Kertarajasa Jayawardhana.
	\item \textbf{Puncak Kejayaan:} Tercapai pada abad ke-14 M di bawah pemerintahan Raja Hayam Wuruk (bergelar Sri Rajasanagara) didampingi oleh Mahapatih Gajah Mada.
	\item \textbf{Sumpah Palapa:} Gajah Mada mengucapkan Sumpah Palapa, sebuah ikrar politik ikonik untuk tidak memakan buah palapa (menikmati kesenangan duniawi) sebelum berhasil menyatukan seluruh wilayah Nusantara di bawah naungan panji Majapahit.
	\item \textbf{Sumber Sejarah:} Kitab \textit{Negarakertagama} (karya Empu Prapanca) yang memuat struktur wilayah dan pemerintahan Majapahit, Kitab \textit{Sutasoma} (karya Empu Tantular) yang memuat frasa toleransi \textit{Bhinneka Tunggal Ika}, serta situs arkeologi luas di Trowulan.
	\item \textbf{Sistem Ekonomi:} Bersifat agraris maritim, menggabungkan hasil bumi padi pedalaman Jawa dengan kendali penuh atas bandar-bandar niaga di pesisir utara Jawa.
\end{itemize}

\section{Masuknya Islam Ke Indonesia (Secara Perlahan)}
Agama Islam mudah diterima oleh masyarakat Nusantara karena beberapa faktor sosiologis dan politik:
\begin{itemize}
	\item Syarat masuk agama Islam sangat mudah, yaitu hanya dengan mengucapkan kalimat Syahadat.
	\item Upacara keagamaan dan ibadahnya tergolong sangat sederhana serta tidak mengenal pembagian kasta sosial.
	\item Islam menyebar secara damai dengan melakukan akulturasi dengan adat, kesenian, dan tradisi lokal yang sudah ada sebelumnya.
	\item Faktor politik memperlancar penyebaran seiring dengan kemunduran dan runtuhnya pengaruh dua kerajaan besar, Majapahit dan Sriwijaya.
\end{itemize}

\subsection{Teori Masuknya Islam}
\begin{itemize}
	\item \textbf{Teori Gujarat:} Islam masuk ke Indonesia dibawa oleh para pedagang Muslim dari Gujarat, India pada abad ke-13 Masehi.
	\item \textbf{Teori Persia:} Islam masuk pada abad ke-13 Masehi dari Persia (Iran) karena adanya kesamaan tradisi keagamaan, seperti perayaan Syura atau 10 Muharram.
	\item \textbf{Teori Makkah (Murni/Arab):} Islam masuk langsung dari Arab (Makkah/Madinah) sejak abad ke-7 Masehi (abad ke-1 Hijriah) dibawa oleh para musafir dan pedagang Arab jalur laut.
	\item \textbf{Teori Cina:} Islam menyebar melalui jalur migrasi etnis Cina Muslim dari daratan Tiongkok ke wilayah pesisir Nusantara sejak masa Dinasti Tang.
\end{itemize}

\subsection{Bukti Awal Penyebaran}
\begin{itemize}
	\item Catatan Dinasti Tang (abad ke-7 M) yang menyebutkan adanya pemukiman pedagang Arab Muslim (*Ta-shih*) di Baros, pantai barat Sumatra.
	\item Catatan perjalanan Marcopolo (1292 M) yang singgah di Kerajaan Perlak dan menjumpai masyarakat setempat yang sudah memeluk Islam.
	\item Catatan perjalanan Ma Huan (musafir utusan Laksamana Cheng Ho) mengenai keberadaan komunitas Muslim di pesisir Jawa.
	\item Kitab \textit{Suma Oriental} karya Tome Pires yang mengisahkan peta penyebaran Islam di jalur niaga Nusantara.
	\item Tulisan huruf Arab kuno pada batu nisan Fatimah binti Maimun di Leran, Gresik (bertahun 1082 M).
	\item Kompleks Makam Muslim kuno di Tralaya dan Trowulan yang berada di dalam lingkungan istana Majapahit.
\end{itemize}

\subsection{Saluran Penyiaran Islam}
\begin{itemize}
	\item \textbf{Perdagangan:} Melalui interaksi antara pedagang asing Muslim dengan penduduk lokal di kota-kota pelabuhan pesisir pantai.
	\item \textbf{Perkawinan:} Pernikahan antara pedagang/ulama Muslim kaya dengan putri bangsawan atau raja lokal.
	\item \textbf{Pendidikan:} Melalui pendirian lembaga pendidikan pesantren, surau, atau dayah untuk mencetak kader ulama baru.
	\item \textbf{Kesenian:} Dakwah budaya menggunakan instrumen gamelan, pertunjukan wayang kulit oleh Sunan Kalijaga, seni sastra suluk, dan tembang makpat.
	\item \textbf{Politik:} Ketika seorang raja atau adipati memeluk Islam, rakyatnya akan mengikuti secara massal demi legitimasi kekuasaan.
\end{itemize}

\subsection{Pemegang Peran Penyebaran Agama Islam}
\subsubsection{Ulama dan Wali Songo}
Dakwah di Nusantara digerakkan oleh para ulama penyiar agama Islam dari Arab, Persia, dan India. Di Pulau Jawa, peran dakwah didominasi oleh sembilan mubaligh agung yang disebut Wali Songo. Setiap wali memiliki strategi dakwah kultural tersendiri untuk melanjutkan penyebaran Islam ke pedalaman jaring-jaring budaya Jawa.
\subsubsection{Perdagangan Internasional}
Kehidupan para pedagang Muslim diterima dengan baik oleh penguasa lokal karena tidak ada pemisahan kaku antara aktivitas ekonomi niaga dengan kewajiban syiar agama untuk menyampaikan ajaran Islam kepada sesama manusia.
\subsubsection{Komunitas Muslim Cina}
Etnis Cina Muslim banyak mewarnai kehidupan sosial-politik publik di masyarakat pesisir Jawa. Komunitas etnis Cina terjalin erat melalui kota pelabuhan ekspor di Guangzhou (Kanton) yang terhubung ke Nusantara melalui ikatan pernikahan campuran dengan penduduk pribumi.

\section{Sejarah Kerajaan Islam di Nusantara}
\subsection{Kesultanan Samudera Pasai}
Kesultanan Samudera Pasai merupakan kerajaan Islam pertama di Nusantara yang berdiri sekitar abad ke-13 Masehi (1267 M). Kesultanan ini terletak di pesisir utara Pulau Sumatra, tepatnya di sekitar Lhokseumawe, Aceh.
\begin{itemize}
	\item \textbf{Pendiri:} Marah Silu, yang setelah memeluk agama Islam berkat bimbingan Syekh Ismail (utusan dari Syarif Makkah), naik takhta dengan gelar resmi Sultan Malik As-Saleh.
	\item \textbf{Kehidupan Ekonomi:} Menjadi salah satu bandar transito komersial internasional terpenting di Selat Malaka. Komoditas ekspor utamanya adalah lada. Pasai juga menjadi kerajaan pertama di Nusantara yang mencetak mata uang emasnya sendiri yang disebut \textit{Dirham}.
	\item \textbf{Sosial-Keagamaan:} Menganut aliran Mazhab Syafi'i. Catatan penjelajah terkenal asal Maroko, Ibnu Battutah, menyebutkan bahwa Sultan Pasai adalah sosok yang sangat saleh, rendah hati, dan suka berdiskusi ilmiah mengenai hukum-hukum Islam dengan kaum ulama.
\end{itemize}

\subsection{Kesultanan Demak}
Kesultanan Demak adalah kerajaan bercorak Islam pertama di Pulau Jawa yang berdiri pada akhir abad ke-15 Masehi (sekitar tahun 1478 atau 1500 M).
\begin{itemize}
	\item \textbf{Pendirian:} Didirikan oleh Raden Patah, yang merupakan putra dari raja Majapahit terakhir (Brawijaya V). Pendirian Demak didukung penuh oleh para bupati pesisir utara Jawa yang telah memeluk Islam serta para Wali Songo.
	\item \textbf{Puncak Kejayaan:} Tercapai pada masa pemerintahan Sultan Trenggono. Di bawah kepemimpinannya, Demak berhasil membendung pengaruh kolonial Portugis dengan mengirimkan armada militer di bawah panglima Fatahillah untuk merebut Sunda Kelapa pada tahun 1527 (kemudian diubah namanya menjadi Jayakarta).
	\item \textbf{Peran Budaya:} Masjid Agung Demak didirikan sebagai pusat ibadah sekaligus markas koordinasi dakwah kultural Islam oleh Wali Songo di seluruh penjuru tanah Jawa.
\end{itemize}

\subsection{Kesultanan Mataram Islam}
Kesultanan Mataram Islam berdiri pada akhir abad ke-16 Masehi (1588 M) yang berpusat di Kotagede, Yogyakarta. Kerajaan ini diawali dari pembukaan Alas Mentaok oleh Ki Ageng Pemanahan sebagai hadiah dari Kesultanan Pajang.
\begin{itemize}
	\item \textbf{Pendiri:} Raden Mas Sutawijaya, putra Ki Ageng Pemanahan, yang naik takhta dengan gelar Panembahan Senopati.
	\item \textbf{Puncak Kejayaan:} Dicapai di bawah kepemimpinan Sultan Agung Hanyokrokusumo (1613--1645 M). Beliau berhasil menyatukan hampir seluruh tanah Jawa, menciptakan Kalender Jawa (perpaduan kalender Saka dan Hijriah), serta memimpin serangan militer kolosal sebanyak dua kali ke Batavia (tahun 1628 dan 1629) guna mengusir kongsi dagang VOC Belanda.
	\item \textbf{Keruntuhan akibat Kolonialisme:} Akibat politik adu domba (*devide et impera*) VOC, Mataram Islam akhirnya pecah melalui Perjanjian Giyanti pada tahun 1755 menjadi dua entitas terpisah: Kasunanan Surakarta dan Kasultanan Yogyakarta.
\end{itemize}

\subsection{Kesultanan Gowa-Tallo}
Kesultanan Gowa-Tallo, yang sering disebut sebagai Kesultanan Makassar, merupakan kerajaan maritim bercorak Islam terbesar di wilayah Indonesia Timur yang mengalami perkembangan pesat pada abad ke-17 Masehi.
\begin{itemize}
	\item \textbf{Karakteristik:} Merupakan hasil penggabungan dua kerajaan kembar (Gowa dan Tallo) yang bersepakat untuk bersama-sama mengelola roda pemerintahan. Raja Gowa, Daeng Manrabia, menjadi Sultan Makassar pertama dengan gelar Sultan Alauddin.
	\item \textbf{Raja Terkenal:} Sultan Hasanuddin (memerintah 1653--1669 M). Karena keberanian, ketangguhan, dan penolakannya terhadap monopoli perdagangan rempah-rempah Belanda, ia dijuluki oleh pihak VOC sebagai \textit{De Haantjes van het Oosten} (Ayam Jantan dari Timur).
	\item \textbf{Kehidupan Ekonomi:} Berfungsi sebagai bandar niaga internasional bebas (*free port*) yang sangat ramai di Sulawesi Selatan, menjembatani jalur pelayaran rempah-rempah dari Maluku menuju kawasan Indonesia Barat dan dunia luar. Namun, dominasi maritim ini patah setelah VOC bersekutu dengan Arung Palakka dari Bone, memaksa Hasanuddin menandatangani Perjanjian Bongaya pada tahun 1667.
\end{itemize}

\end{document}